\section{Introduction}
\label{sec:intro}

A reader who learns that ``peanuts are legumes'' or ``Georgia is both a US state and a country'' rarely thinks of these as deep facts: they are textbook standard knowledge, the cognitive currency of an everyday conversation. Large language models clearly contain this knowledge --- a 2.8B parameter model answers head-of-government, capital-of, and classification-level queries with non-trivial accuracy. But \emph{how much} of the model's stored capacity is devoted to it? Is the standard-knowledge core a substantial fraction of the bit budget, or a negligible footnote dwarfed by long-tail facts and linguistic memorization?

\para{Why the question matters.} The answer changes how we reason about distillation, quantization, and scaling. If standard knowledge consumes \(10\%\) of the bit budget, aggressive pruning or \(\text{int}4\) quantization may damage the everyday-fact core that users notice most. If it consumes a vanishing fraction, that core is essentially free, and the parameter pressure of scaling comes from somewhere else entirely. The question also frames how to interpret existing benchmarks: an \(8.9\%\) \bear accuracy at \(410\)M parameters is striking when contextualized as ``\(8.9\%\) of \(7{,}731\) facts'', but it tells us nothing about the parameter footprint until the bits are converted to a fraction of capacity.

\para{The gap.} Three rapidly converging literatures touch this question without answering it directly. Capacity scaling laws~\citep{allenzhu2024knowledge,morris2025memorize} measure raw bits per parameter --- \(2\) bpp for retrievable knowledge tuples, \(3.6\) bpp for raw memorization --- but never partition stored bits by knowledge type. Mechanistic interpretability~\citep{geva2021ffnkv,dai2022knowledge,meng2022rome,hong2025specialization} localizes \emph{individual} facts to MLP layers and FFN value vectors but stops at the per-fact level. Probing benchmarks~\citep{petroni2019lama,wiland2024bear,powell2024taxi} report what fraction of facts the model knows but never convert that into a parameter footprint. The missing measurement is the bridge between the two: \emph{what fraction of an LLM's parameters is devoted to standard categorical and conditional knowledge?}

\para{Our approach.} We combine log-likelihood probing with the \(2\)-bpp ceiling to do bit-budget accounting (\Figref{fig:scaling_bits}). For each model and each (subject, relation, gold answer) tuple in \bear (relational knowledge) and \taxi (categorical knowledge), we score every candidate completion under the relation template and convert above-chance accuracy into \emph{acquired bits}. Summing across all relations and dividing by \(2 \cdot |\text{params}|\) gives the parameter share. We run the analysis across the \pythia model family --- \(160\)M, \(410\)M, \(1\)B, \(2.8\)B --- so that we can also report how the share scales.

\para{Headline result.} Standard categorical and relational knowledge consumes only \(\mathbf{0.00006\%}\) to \(\mathbf{0.00022\%}\) of the bit budget at \(2\) bpp --- \emph{roughly one part in a million of stored capacity}. The share decreases monotonically with parameter count from \(410\)M onward; bootstrap \(95\%\) confidence intervals at \(1\)B and \(2.8\)B are disjoint from those at \(160\)M and \(410\)M, so the decay is statistically reliable despite only four model points. The mechanism is sub-linear scaling of acquired bits: a log-log regression on \bear bits gives a slope of \(0.54\) (\(p=0.04\)) --- adding a parameter buys only half a fractional bit of standard knowledge, far less than the one bit of capacity it adds.

\para{Contributions.} In summary:
\begin{itemize}[leftmargin=*,itemsep=0pt,topsep=0pt]
    \item We propose a probe-based bit-budget accounting that converts \bear and \taxi log-likelihood probe accuracy into a fraction of an LLM's stored capacity, using the \(2\)-bpp ceiling of \citet{allenzhu2024knowledge}.
    \item We conduct the analysis across four \pythia model scales (\(160\)M, \(410\)M, \(1\)B, \(2.8\)B), reporting bootstrap \(95\%\) confidence intervals and a log-log scaling fit.
    \item We complement the bit-budget figure with an analytical per-fact footprint cross-check derived from \kn~\citep{dai2022knowledge} and \rome~\citep{meng2022rome}, showing that those upper bounds over-estimate the aggregate footprint by \(\mathbf{4}\) orders of magnitude --- evidence that parameter sharing across facts is the dominant regime in modern LLMs.
    \item We release the probing code, per-relation breakdowns, and bootstrap CIs at the project repository, so all numbers in the paper are directly reproducible.
\end{itemize}

\para{Organization.} \Secref{sec:related} positions our work against capacity-scaling, mechanistic, and probing literatures. \Secref{sec:method} formalizes the bit-budget accounting and lists models, datasets, and metrics. \Secref{sec:results} reports the headline scaling table, regression fit, and per-relation breakdowns. \Secref{sec:discussion} interprets the result, lists six limitations, and outlines what occupies the rest of the bit budget. \Secref{sec:conclusion} summarizes and points to follow-ups across architectures, templates, and quantization regimes.
